\section{Harmonic fillings and mixed state correlations}
\label{hfill:section}

The two states in Chapter~\ref{hst:chapter} have identical
one-mode density operators, but a mixed insertion distinguishes them.
To interpret that insertion geometrically, we need a space of fillings,
a measure, and an action of functions on its state space.
A real scalar field on the disk supplies these data.
Its harmonic extension realizes the mixed insertion as a function
of two jets at the marked point.

There is a separate question at a collision.
The boson--fermion map extends to principal values and their contacts
on charts with independent normal coordinates.
The contact module, however, excludes a positive-definite pairing with the
pointwise coordinate involution.
We first construct this relative comparison, then construct the
scalar filling and its positive sewing.

\subsection{The state map and relative contacts}

Work at level one and charge zero over $\mathbb C$.
Let $\mathcal P=\mathbb C[p_1,p_2,\ldots]$, with
$\deg p_n=n$ and Hall form
\[
 \langle p^{\boldsymbol m},p^{\boldsymbol l}\rangle
 =\delta_{\boldsymbol m,\boldsymbol l}
       \prod_{n\geq1}n^{m_n}m_n!.
\]
The form is conjugate linear in its first argument.
Write $\mathcal H_B$ for its Hilbert completion and $E$ for energy.
On $\mathcal P$ set
\[
 a_{-n}=p_n,\qquad a_n=n\partial_{p_n}\quad(n>0),
 \qquad a_0=0,
 \qquad T_B=\sum_{n\geq1}n p_{n+1}\partial_{p_n}.
\]

Let $\mathcal F_0$ be the charge-zero semi-infinite wedge, with
orthonormal basis indexed by partitions $\lambda$:
\[
 |\lambda\rangle_F
 =e_{\lambda_1-1}\wedge e_{\lambda_2-2}\wedge\cdots.
\]
Its tail is $e_{-i}$ and its Hilbert completion is $\mathcal H_F$.
On the sum of charge sectors, $\psi_r^\dagger$ inserts $e_r$ and
$\psi_r$ removes it, with the exterior signs.
For $n\ne0$ put
\[
 A_n=\sum_{r\in\mathbb Z}\psi_{r-n}^\dagger\psi_r,
 \qquad
 T_F=\sum_{r\in\mathbb Z}(r+1)\psi_{r+1}^\dagger\psi_r.
\]
These sums are finite on each finite-energy vector.
The zero current acts by zero on this charge sector.

The Hall--wedge map of \eqref{hst:eq:fermion-map} is
\begin{equation}\label{hfill:eq:U}
 U:\mathcal P\longrightarrow\mathcal F_0,
 \qquad U(s_\lambda)=|\lambda\rangle_F.
\end{equation}
It preserves the vacuum and intertwines $a_n$ with $A_n$.
It also intertwines $T_B$ with $T_F$.
Indeed both translations kill the vacuum and satisfy
$[T,a_n]=-n a_{n-1}$, with the corresponding fermionic formula.
Commuting through the creation operators proves the assertion.

The map extends to a unitary $\mathcal H_B\to\mathcal H_F$.
Finite occupation vectors are a common core for each closed oscillator.
The finite-energy intertwining therefore identifies their closed graphs
and the domains of their successive products.
This retains the operator domains specified in Chapter~\ref{hst:chapter}.

Fix a polydisk with independent normal coordinates
$u_1,\ldots,u_d$ and additional coordinates $w$.
Let $\mathcal O$ denote its jointly holomorphic functions.
For one normal define the two-term complex
\begin{equation}\label{hfill:eq:contact-complex}
 \mathcal E_i^0=\mathcal O[u_i^{-1}],\qquad
 \mathcal E_i^1=\mathcal O[u_i^{-1}]/\mathcal O,
 \qquad d f=f\bmod\mathcal O.
\end{equation}
The degrees are $0$ and $1$.

Realize its first term by principal-value currents and put
\begin{equation}\label{hfill:eq:residue}
 r_k(u)=\bar\partial\operatorname{PV}(u^{-k})
 =\frac{\pi(-1)^{k-1}}{(k-1)!}
       \partial_u^{k-1}\delta^{(2)}(u)\,d\bar u.
\end{equation}
The class of $u^{-k}$ in the second term maps to $r_k$.
Here $\bar\partial\operatorname{PV}(1/u)
=\pi\delta^{(2)}(u)d\bar u$ fixes the normalization.
Holomorphic multiplication and differentiation respect this realization,
by \eqref{scmp:eq:jets} and its residue-current realization.

Set
\[
 \mathcal E=\mathcal E_1\otimes_{\mathcal O}\cdots
                         \otimes_{\mathcal O}\mathcal E_d,
 \qquad
 C_B=\mathcal E\otimes_{\mathbb C}\mathcal P,
 \qquad C_F=\mathcal E\otimes_{\mathbb C}\mathcal F_0.
\]
The tensor differential has the cohomological sign.
For two normals, for example,
\[
 d(\operatorname{PV}(u^{-a})\operatorname{PV}(v^{-b}))
 =r_a(u)\operatorname{PV}(v^{-b})
   +\operatorname{PV}(u^{-a})r_b(v).
\]
Applying $d$ again gives opposite corner terms.
Products here involve independent normal coordinates.
The same-normal meromorphic multiplication on contacts is excluded by
Proposition~\ref{scmp:prop:no-product}.

Every section has finite pole order and finite state support.
On each such subspace use the compact-open topology of its holomorphic
coefficients, then take the locally convex inductive limit.
The coefficient currents act on compactly supported test forms.
We retain each face submodule, supported on a specified intersection
of the divisors $u_i=0$.
The filtration by at least $j$ contact factors is a subcomplex filtration.

\begin{proposition}\label{hfill:prop:relative}
The map
\[
 \widetilde U=\operatorname{id}_{\mathcal E}\otimes U:
 C_B\longrightarrow C_F
\]
is a continuous isomorphism of cochain $\mathcal O$-modules.
It preserves the face submodules, contact filtration, current modes,
holomorphic translations, state translation, and vacuum.
Its tensor powers preserve external products in independent normal charts.
The weak finite-energy matrix coefficients of ordered current fields
and polynomial descendants have matching principal values and contacts.
\end{proposition}

\begin{proof}
The coefficient differential and its contact signs act on $\mathcal E$.
The map $U$ acts in degree zero on the finite-dimensional energy spaces.
Thus it commutes with the differential, coefficient operations, and
the topology just specified.
Equation~\eqref{hfill:eq:U} gives the state operations and their domains.
Tensoring these identities gives the external-product assertion.

Write $J_B(z)=\sum_n a_nz^{-n-1}$ and define $J_F$ using $A_n$.
The mode identity gives $UJ_B(z)=J_F(z)U$ coefficientwise.
Moving annihilators to the right gives contraction $(z-w)^{-2}$ on
both sides.
For polynomial vectors, the remaining matrix coefficients are finite
Laurent polynomials after this contraction is extracted.

Wick expansion and differentiation therefore identify the rational
matrix coefficients of every finite ordered product of descendants.
Their expansions use the same region $|z_1|>\cdots>|z_m|$.
On an independent-normal chart, extend the common coefficient by
\eqref{hfill:eq:contact-complex}.
The principal values and all contact differentials then agree.
\end{proof}

This comparison specifies coefficient cochains and weak field coefficients.
It supplies no multiplication of arbitrary singular fields by supported
currents in the same normal coordinate.
An independently defined native quantum cochain complex requires its
own comparison with $C_B$ or $C_F$.

\subsection{A pairing obstruction on the marked residue module}

In one normal let $\mathcal R$ be the span of the currents $r_k$.
Ordinary multiplication by a smooth function acts on these representatives.
In particular,
\begin{equation}\label{hfill:eq:marked-actions}
 u r_1=0,\qquad u r_k=r_{k-1}\ (k>1),
 \qquad \bar u r_k=0.
\end{equation}
These identities use the pointwise involution $u^*=\bar u$.

For the support calculation, let $\mathcal J^0$ be the finite span
of the jets $\partial_u^a\partial_{\bar u}^b\delta^{(2)}(u)$,
and put $\mathcal J^1=\mathcal J^0d\bar u$.
Its differential raises $b$ by one.
Consequently
\[
 H^1(\mathcal J)=
 \bigoplus_{a\geq0}\mathbb C
 [\partial_u^a\delta^{(2)}(u)d\bar u].
\]
Evaluation on a holomorphic test jet detects each indicated class.
In particular $[r_1]\ne0$.
The same current is exact in the larger principal-value complex,
where $d\operatorname{PV}(1/u)=r_1$.

\begin{proposition}\label{hfill:prop:positive-obstruction}
There is no positive-definite Hermitian form on $\mathcal R$
for which $M_u^*=M_{\bar u}$.
The same obstruction holds in any marked module retaining
\eqref{hfill:eq:marked-actions} and a nonzero image of $r_1$.
In particular it is preserved by module comparisons that intertwine
these coordinate actions and retain the detected residue class.
\end{proposition}

\begin{proof}
The required adjoint identity gives
\[
 \|r_1\|^2=\langle ur_2,ur_2\rangle
 =\langle r_2,\bar u u r_2\rangle=0.
\]
The residue detector shows $r_1\ne0$.
The same calculation applies to its nonzero marked image.
\end{proof}

The proposition concerns the marked module of representatives.
Smooth multiplication does not act on general Dolbeault cohomology:
$\bar\partial M_{\bar u}-M_{\bar u}\bar\partial$
is multiplication by $d\bar u$.
Thus the supported cohomology calculation cannot supply that action.
Reflected pairings between different insertion domains remain possible,
as in Section~\ref{hss:sec:pairing}.

\subsection{The scalar filling and its state space}

Let $D$ be the unit disk with its center and angular coordinate marked.
Fix the constant mode to zero.
For finitely many complex numbers $Z_n$, define the real harmonic field
\begin{equation}\label{hfill:eq:harmonic}
 \phi=h+\bar h,\qquad
 h(z)=\sum_{n\geq1}\frac{Z_n}{n}z^n,
 \qquad
 S_D(\phi)=\frac1{8\pi}\int_D|\nabla\phi|^2\,d^2z.
\end{equation}
The metric is Euclidean and $d^2z=dx\,dy$.
The action determines the ground amplitude $\exp(-S_D)$.
At a finite cutoff, normalize its squared absolute value against
Lebesgue measure in the coordinates $Z_n$.

\begin{proposition}\label{hfill:prop:measure}
These finite laws are consistent and define
\begin{equation}\label{hfill:eq:measure}
 d\mu=\bigotimes_{n\geq1}
       \frac1{\pi n}e^{-|Z_n|^2/n}\,d^2Z_n.
\end{equation}
Almost every sequence determines a smooth harmonic interior and a
distributional boundary value through \eqref{hfill:eq:harmonic}.
Almost every such field has infinite Dirichlet energy.
\end{proposition}

\begin{proof}
Since $|\nabla\phi|^2=4|h'|^2$, angular orthogonality gives
\[
 S_D(\phi)=\frac12\sum_n\frac{|Z_n|^2}{n}.
\]
Squaring the amplitude gives \eqref{hfill:eq:measure} at every cutoff.
Each unused factor has mass one, which proves consistency.

The Gaussian tail is
$\mu(|Z_n|>n)=e^{-n}$.
Hence $|Z_n|\leq n$ eventually, almost surely.
The power series and all its derivatives converge uniformly on
every smaller disk.
Its boundary Fourier coefficients have polynomial growth, so they
define a distribution.

The variables $|Z_n|^2/n$ are independent exponentials of mean one.
Their partial sums diverge almost surely by the strong law.
The nonnegative Fourier expansion of the Dirichlet integral then
gives infinite total energy.
\end{proof}

Let $\mathcal H_{\mathrm{hol}}$ be the closed subspace of
$L^2(\mu)$ generated by holomorphic cylinder polynomials.
This defines the recipient independently of Hall states.
Gaussian polar integration gives
\[
 \int\overline{Z^{\boldsymbol a}}Z^{\boldsymbol b}\,d\mu
 =\delta_{\boldsymbol a,\boldsymbol b}
       \prod_n n^{a_n}a_n!.
\]
Therefore substitution extends uniquely to a unitary
\begin{equation}\label{hfill:eq:V}
 V:\mathcal H_B\longrightarrow\mathcal H_{\mathrm{hol}},
 \qquad V(p_n)=Z_n.
\end{equation}
Under this map the filling energy is the closure of
$E=\sum_n nZ_n\partial_{Z_n}$ on cylinder polynomials.

\subsection{Geometric insertions and the mixed correlation}

Let $P_{\mathrm{hol}}$ be orthogonal projection onto
$\mathcal H_{\mathrm{hol}}$.
For a cylinder polynomial $f(Z,\bar Z)$, define
$T_f=P_{\mathrm{hol}}M_f$ on holomorphic cylinder polynomials.
Multiplication takes this domain into $L^2(\mu)$.
The projection is therefore defined there.

\begin{proposition}\label{hfill:prop:toeplitz}
For finite multi-indices $\alpha,\beta$,
\[
 T_{Z^\alpha\bar Z^\beta}
 =\prod_n(n\partial_{Z_n})^{\beta_n}\,M_{Z^\alpha}.
\]
For cylinder polynomial symbols, $T_fT_g=T_{f\star g}$, where
\begin{equation}\label{hfill:eq:star}
 f\star g=
 \sum_\alpha\frac{(-1)^{|\alpha|}\prod_n n^{\alpha_n}}{\alpha!}
       (\partial_Z^\alpha f)(\partial_{\bar Z}^\alpha g).
\end{equation}
The sum is finite.
In particular $Z_n\star\bar Z_n=|Z_n|^2-n$.
\end{proposition}

\begin{proof}
Pair a projected monomial against every holomorphic monomial.
The Gaussian moment formula gives
$P_{\mathrm{hol}}(\bar Z_n^b F)
=(n\partial_{Z_n})^bF$ for holomorphic polynomial $F$.
This proves the first formula.

In this ordering, derivatives stand to the left of multiplications.
To compose two such expressions use
$M_{Z_n}(n\partial_{Z_n})
=(n\partial_{Z_n})M_{Z_n}-n$.
The finite binomial expansion for each coordinate is
\eqref{hfill:eq:star}.
Different coordinates commute, which proves the multivariable identity.
\end{proof}

Put
\[
 u_\pm=\frac{p_1^2\pm p_2}{2},\qquad
 F_\pm=Vu_\pm=\frac{Z_1^2\pm Z_2}{2},\qquad
 d\nu_\pm=|F_\pm|^2d\mu.
\]
Both states have norm one.
The marked filling has the real observable
\begin{equation}\label{hfill:eq:mixed}
 \mathcal C(h)=2\operatorname{Re}
       \bigl(h'(0)^2\overline{h''(0)}\bigr)
 =Z_1^2\bar Z_2+\bar Z_1^{\,2}Z_2.
\end{equation}
It is invariant under $z\mapsto e^{i\theta}z$.
It depends on relative mode phases, which individual mode laws can lose.

\begin{proposition}\label{hfill:prop:mixed}
On the polynomial domain,
\[
 T_{\mathcal C}=VOV^{-1},\qquad
 O=a_{-1}^{\,2}a_2+a_{-2}a_1^{\,2}.
\]
The two filling laws satisfy
\begin{equation}\label{hfill:eq:expectations}
 \int\mathcal C\,d\nu_+=2,\qquad
 \int\mathcal C\,d\nu_-=-2.
\end{equation}
Their single-coordinate marginals agree.
For coordinates $1$ and $2$, their densities against $\mu_1$ and
$\mu_2$ are respectively
\[
 \frac{|Z_1|^4+2}{4},\qquad
 \frac{|Z_2|^2+2}{4}.
\]
Every other coordinate has its vacuum Gaussian law.
\end{proposition}

\begin{proof}
The two different modes commute in the Toeplitz formula, giving
the stated operator.
Write $A=Z_1^2$ and $B=Z_2$.
Then $\mathbb E|A|^2=\mathbb E|B|^2=2$ and
$\mathbb E(A\bar B)=0$.
This proves the normalization and both marginal formulas.

The only surviving terms in
$\mathbb E[(A\bar B+\bar A B)|A\pm B|^2]/4$
are $\pm2\mathbb E(|A|^2|B|^2)/4=\pm2$.
Replacing the plus in \eqref{hfill:eq:mixed} by a minus gives zero
on both laws, so it fails to represent the fixed Hermitian operator $O$.
\end{proof}

\subsection{Recovering a density operator from its filling law}

The unitary map $V$ compares amplitudes.
A mixed state instead determines a probability law.
For a positive trace-class operator on $\mathcal H_{\mathrm{hol}}$,
choose a spectral decomposition
\[
 \rho=\sum_j\lambda_j|F_j\rangle\langle F_j|,
 \qquad \lambda_j\geq0,
 \qquad \|F_j\|=1,
\]
and define
\begin{equation}\label{hfill:eq:density-law}
 D_\rho=\sum_j\lambda_j|F_j|^2,
 \qquad d\nu_\rho=D_\rho\,d\mu.
\end{equation}
The sum converges in $L^1(\mu)$ by positivity.
For a bounded measurable function $f$, compression of $M_f$ is bounded
and
\[
 \int f\,d\nu_\rho
 =\operatorname{Tr}(\rho P_{\mathrm{hol}}M_f|_{\mathcal H_{\mathrm{hol}}}).
\]
This identity makes the measure independent of the spectral decomposition.

\begin{theorem}\label{hfill:thm:density-recovery}
The map $\rho\mapsto\nu_\rho$ is injective on positive trace-class
operators.
It is additive and positively homogeneous, and
$\nu_\rho(\mathbb C^{\mathbb N})=\operatorname{Tr}\rho$.
In particular it is an affine injection from density operators to
probability laws of harmonic fillings.
\end{theorem}

\begin{proof}
First retain only the first $N$ modes.
Their holomorphic space has reproducing vector
\[
 K_z(w)=\exp\left(\sum_{n=1}^N\frac{\bar z_nw_n}{n}\right),
 \qquad
 \|K_z\|^2=\exp\left(\sum_{n=1}^N\frac{|z_n|^2}{n}\right).
\]
For a trace-class operator $A$, its diagonal
$\langle K_z,AK_z\rangle$ is real analytic.
Its series converges locally uniformly by the Hilbert--Schmidt
inequality applied to the matrix of $A$ and $K_z\otimes\overline{K_z}$.

In the normalized monomial basis, its coefficient of
$z^\alpha\bar z^\beta$ is
\[
 \frac{A_{\alpha\beta}}
 {\sqrt{\prod_n n^{\alpha_n+\beta_n}\alpha_n!\beta_n!}}.
\]
Wirtinger derivatives at zero recover every matrix entry.
Equality almost everywhere of two such diagonals implies equality
everywhere, by continuity and the full support of the Gaussian measure.
Thus the finite-mode law determines the finite-mode operator.

The marginal of \eqref{hfill:eq:density-law} on these $N$ coordinates
is the diagonal of the reduced operator $\rho_N$.
To see this, test against a bounded function of those coordinates.
The compressed multiplication is its finite-mode compression tensored
with the identity on the remaining modes.
The defining property of the partial trace gives the assertion.

It remains to recover the full operator from all $\rho_N$.
Let $\alpha,\beta$ be finite occupation sequences.
For $N$ containing their supports, form the bounded local matrix unit
\[
 L_N=|\alpha_{\leq N}\rangle\langle\beta_{\leq N}|
                         \otimes1_{>N}.
\]
It includes the vacuum occupations up to $N$.
On every finite occupation vector it eventually agrees with
$|\alpha\rangle\langle\beta|$.
Since $\|L_N\|=1$, this is strong operator convergence.

Pairing a uniformly bounded strongly convergent sequence with a
trace-class operator preserves the limit.
This follows by finite-rank approximation in trace norm.
Hence the reduced operators determine every global occupation matrix entry.
They determine $\rho$.
Additivity, homogeneity and the mass identity follow directly from
the bounded-function identity preceding the theorem.
\end{proof}

\subsection{Scalar-cylinder sewing}

The disk action also has an independent cylindrical transfer.
Use coordinates $(t,\theta)$ with $\theta$ of period $2\pi$ and
flat conformal metric $dt^2+d\theta^2$.
At a finite mode cutoff write
\[
 \phi(t,\theta)=\sum_n
 \left(\frac{Z_n(t)}n e^{in\theta}
       +\frac{\overline{Z_n(t)}}n e^{-in\theta}\right).
\]
The real scalar action is
\begin{equation}\label{hfill:eq:cylinder-action}
 S_{\mathrm{cyl}}
 =\frac1{8\pi}\int
       (|\partial_t\phi|^2+|\partial_\theta\phi|^2)
                                      \,dt\,d\theta
 =\frac12\sum_n\int
       \left(\frac{|\dot Z_n|^2}{n^2}+|Z_n|^2\right)dt.
\end{equation}
Thus each complex mode is a two-dimensional real oscillator of
mass $n^{-2}$ and frequency $n$.

Its Schrödinger Hamiltonian and normalized ground state are
\[
 H_n=-\frac{n^2}{2}(\partial_x^2+\partial_y^2)
                   +\frac{x^2+y^2}{2},
 \qquad
 \psi_n(z)=\frac1{\sqrt{\pi n}}e^{-|z|^2/(2n)},
 \qquad H_n\psi_n=n\psi_n.
\]
The positive self-adjoint realization is the oscillator closure
from the Schwartz domain.
Normalize its transfer by this ground energy at every finite cutoff.

For $0<|q|<1$, let $\eta_n$ be independent circular complex Gaussians
of variance one, independent of $Z$.
Define on bounded cylinder functions
\begin{equation}\label{hfill:eq:mehler}
 (P_qF)(Z)=\mathbb E_\eta
 F\left(\left(q^nZ_n+\sqrt{n(1-|q|^{2n})}\,\eta_n
                                                   \right)_{n\geq1}\right).
\end{equation}
The phase of $q$ records angular rotation.
For real $q=e^{-t}$, the positive parameter $t$ is collar length.

\begin{theorem}\label{hfill:thm:sewing}
The normalized scalar-cylinder transfer is $P_{e^{-t}}$.
The operators in \eqref{hfill:eq:mehler} extend to contractions on
$L^2(\mu)$ and satisfy
\begin{equation}\label{hfill:eq:sewing-laws}
 P_qP_r=P_{qr},\qquad P_q^*=P_{\bar q},\qquad
 P_q|_{\mathcal H_{\mathrm{hol}}}=VE_qV^{-1},
 \quad E_q=q^E.
\end{equation}
For real $0<q<1$, $P_q$ is a positive operator.
The state pairing $B_q(F,G)=\langle F,P_qG\rangle$ then has
\[
 B_q(F_\pm,T_{\mathcal C}F_\pm)=\pm2q^2,
 \qquad B_q(F_\pm,F_\pm)=q^2.
\]
The level-one current $J=J_B$ has covariance
\begin{equation}\label{hfill:eq:covariance}
 \langle1,J(z)E_qJ(w)1\rangle
 =\frac{q}{(z-qw)^2},\qquad |z|>|qw|.
\end{equation}
\end{theorem}

\begin{proof}
Multiplication by $\psi_n$ is a unitary from $L^2(\mu_n)$
to $L^2(\mathbb C,d^2z)$.
Its ground-state transform gives
\[
 \psi_n^{-1}(H_n-n)\psi_n
 =-\frac{n^2}{2}\Delta+n(x\partial_x+y\partial_y).
\]
The Gaussian transition with mean $e^{-nt}z$ and complex variance
$n(1-e^{-2nt})$ solves its heat equation.
Differentiating its density verifies the generator, and the variance
tends to zero at $t=0$.
This proves the finite-cutoff scalar transfer formula.

Every transformed coordinate in \eqref{hfill:eq:mehler} again has law
$\mu_n$, and the coordinates remain independent.
Jensen's inequality therefore gives the $L^2$ contraction.
The cylinder prescriptions are consistent, so they extend by density.
This construction uses normalized finite transfers, without a sum of
infinite vacuum energies.

Composing the two independent noises gives mean $(qr)^nZ_n$ and variance
$n(1-|qr|^{2n})$.
This proves the semigroup law.
Reversing the two coordinates of their joint Gaussian law replaces
$q$ by $\bar q$, which proves the adjoint identity.
Circularity kills every nonconstant holomorphic noise moment.
Thus on a holomorphic monomial the eigenvalue is $q^{\sum_n n m_n}$.

For positivity one can use the complex Hermite polynomials
\[
 H_{a,b}^{(n)}(z,\bar z)
 =\sum_{j=0}^{\min(a,b)}
   (-n)^j j!\binom aj\binom bj z^{a-j}\bar z^{b-j}.
\]
Gaussian integration by parts makes them an orthogonal basis in one mode.
Their eigenvalues under $P_q$ are $q^{n(a+b)}$ for real $q$.
Finite tensor products give a complete basis of $L^2(\mu)$,
with positive eigenvalues when $0<q<1$.

The vectors $F_\pm$ have energy two and
$T_{\mathcal C}F_\pm=\pm2F_\pm$.
This gives the sewn mixed test.
Finally the current commutator and vacuum give
\[
 \langle1,J(z)E_qJ(w)1\rangle
 =\sum_{n\geq1}nq^n z^{-n-1}w^{n-1}
 =\frac{q}{(z-qw)^2}.
\]
\end{proof}

The factor $q$ in \eqref{hfill:eq:covariance} is the one-form factor
in the coordinate pullback $W=qw$.
The corresponding contact lies on $z-qw=0$.
It is compared using that normal coordinate, as in
\eqref{hss:eq:sewn-cov}.
The normalized mixed expectation remains $\pm2$ after sewing.

Propagation of a density operator uses two transfers.
For a density operator $\rho$ on $\mathcal H_B$, set
\begin{equation}\label{hfill:eq:propagated-density}
 \rho_q=
 \frac{q^E\rho\,\bar q^{\,E}}
      {Z_\rho(q)},
 \qquad Z_\rho(q)=\operatorname{Tr}(\rho|q|^{2E}).
\end{equation}
The denominator is positive because $|q|^{2E}$ has no kernel.

\begin{lemma}\label{hfill:lem:collar-moments}
For every $m\geq0$ and fixed $0<|q|<1$,
\[
 \operatorname{Tr}(\rho_q(E+1)^m)
 \leq\frac{\sup_{e\geq0}(e+1)^m|q|^{2e}}{Z_\rho(q)}<\infty.
\]
\end{lemma}

\begin{proof}
Insert the spectral projections of $E$ into the trace.
The nonnegative weights $\operatorname{Tr}(\rho\,1_{E=e})$ sum to one.
The exponential weight bounds each term by the displayed supremum.
\end{proof}

This estimate depends on the collar and on $Z_\rho(q)$.
It does not give a uniform assertion as the collar shrinks.

\subsection{The cubic observable at every energy}

Independently of the filling, the current modes define
\begin{equation}\label{hfill:eq:W}
 W=\frac12\sum_{i,j>0}
       \left(a_{-i}a_{-j}a_{i+j}
                    +a_{-(i+j)}a_i a_j\right).
\end{equation}
It preserves each finite-dimensional energy space.
Write $\mathcal H_e$ for that space and $W_e$ for the restriction.
On $\mathcal H_B$, use the self-adjoint direct sum of these symmetric blocks,
with domain
\[
 \mathcal D(W)=\left\{F=\sum_eF_e\in\mathcal H_B:
                          \sum_e\|W_eF_e\|^2<\infty\right\}.
\]
Its normalization is \eqref{lmte:eq:cut-join}.

On the fermion side define the vacuum-subtracted second sea moment by
\[
 M_2|\lambda\rangle_F
 =\frac12\sum_{i\geq1}
   \left[(\lambda_i-i+\tfrac12)^2-(-i+\tfrac12)^2\right]
                                        |\lambda\rangle_F.
\]
The sum is finite. Use its self-adjoint diagonal extension on $\mathcal H_F$.

\begin{lemma}\label{hfill:lem:cut-join}
One has
\[
 UW=M_2U,\qquad
 Ws_\lambda=c_\lambda s_\lambda,
 \qquad
 c_\lambda=\frac12\sum_i\lambda_i(\lambda_i-2i+1).
\]
Moreover $\|W|_{E=e}\|\leq e(e-1)/2$.
\end{lemma}

\begin{proof}
The eigenvalue agrees with the charges in \eqref{lst:eq:Qcharges}.
It can also be derived directly from finite alternants.
Choose $N\geq e$ and specialize $p_m=\sum_{a=1}^N x_a^m$,
which is injective through energy $e$.
Put $D_a=x_a\partial_{x_a}$ and
$\Delta_N=\prod_{a<b}(x_a-x_b)$.

The chain rule rewrites \eqref{hfill:eq:W} on symmetric polynomials as
\begin{equation}\label{hfill:eq:alternant-operator}
 Wf=\frac1{2\Delta_N}\sum_aD_a^2(\Delta_N f)
      -\frac12\sum_{i=1}^N(N-i)^2f
      -\frac{2N-1}{2}\sum_aD_af.
\end{equation}
To check the first-order terms, use
\[
 \sum_{a<b}\frac{(x_a+x_b)(x_a^m-x_b^m)}{x_a-x_b}
 =(N-m)p_m+\sum_{i=1}^{m-1}p_ip_{m-i}.
\]
The second-order chain rule gives
$\sum_{i,j}ij p_{i+j}\partial_{p_i}\partial_{p_j}$.
Together these are exactly twice the two terms in
\eqref{lmte:eq:cut-join}.

For $f=s_\lambda$, the numerator $\Delta_Nf$ is the alternant
with exponents $\lambda_i+N-i$.
The sum of the $D_a^2$ acts on it by
$\sum_i(\lambda_i+N-i)^2$.
Substitution in \eqref{hfill:eq:alternant-operator} gives $c_\lambda$.
Expanding the two squares in $M_2$ gives the same value.

Finally $c_\lambda=\sum_{(i,j)\in\lambda}(j-i)$.
Its upper bound is
$\sum_i\lambda_i(\lambda_i-1)/2\leq e(e-1)/2$.
Transposing the diagram negates the content sum and gives the lower bound.
The Schur basis is orthonormal, which proves the operator norm bound.
\end{proof}

Deleting the half shift changes $M_2$ by $-E/2$.
It already gives $-1/2$ on the one-box state, where $W$ is zero.
The shift therefore cannot be discarded in this comparison.

For $0<r<1$, define a function of the harmonic filling by
\begin{equation}\label{hfill:eq:radial-observable}
 g_r(\theta)=re^{i\theta}h'(re^{i\theta}),
 \qquad
 \Phi_r=\int_0^{2\pi}
       \operatorname{Re}(g_r^2\overline{g_r})\frac{d\theta}{2\pi}.
\end{equation}
It is real and invariant under rotations of the disk.
Angular integration gives the almost surely convergent expression
\begin{equation}\label{hfill:eq:radial-symbol}
 \Phi_r=\frac12\sum_{i,j\geq1}r^{2(i+j)}
       \left(Z_iZ_j\bar Z_{i+j}
                     +\bar Z_i\bar Z_j Z_{i+j}\right).
\end{equation}
Let $W_r$ be \eqref{hfill:eq:W} with each summand multiplied by
$r^{2(i+j)}$.
Use the same self-adjoint direct-sum construction for $W_r$.

\begin{theorem}\label{hfill:thm:radial-limit}
On finite-energy states, the Toeplitz compression of $\Phi_r$ is
$VW_rV^{-1}$.
For every density operator $\rho$ on $\mathcal H_B$ and every
fixed $0<|q|<1$,
\begin{equation}\label{hfill:eq:radial-limit}
 \lim_{r\uparrow1}\int\Phi_r\,d\nu_{V\rho_qV^{-1}}
 =\operatorname{Tr}(\rho_qW).
\end{equation}
All integrals at $r<1$ are absolutely convergent.
For a finite-energy state no collar is required.
In particular
$\int\Phi_r\,d\nu_\pm=\pm r^4$.
\end{theorem}

\begin{proof}
First truncate the sum in \eqref{hfill:eq:radial-symbol}.
The mode $i+j$ differs from both $i$ and $j$.
Thus the Toeplitz ordering has no extra contraction, and its
compression is the corresponding finite sum in $VW_rV^{-1}$.

For integrability, let $F$ have energy $e$.
Creation in mode $n$ has norm at most $\sqrt{e+n}$ on that energy space.
Put $k=i+j$.
Pointwise absolute values give
\[
 \|Z_iZ_j\bar Z_kF\|_{L^2(\mu)}
 =\|Z_iZ_jZ_kF\|_{L^2(\mu)}
 \leq(e+2k)^{3/2}\|F\|.
\]
The conjugate term has the same bound.

For a general holomorphic vector, the three creation operators shift
each energy by the same amount $2k$.
Their images from different energy spaces are orthogonal.
Consequently
\begin{equation}\label{hfill:eq:multiplication-bound}
 \|\Phi_rF\|_{L^2(\mu)}
 \leq C_r\|(E+1)^{3/2}F\|,
 \qquad
 C_r=\sum_{k\geq2}(k-1)(2k)^{3/2}r^{2k}<\infty.
\end{equation}
This proves convergence of the symbol times $F$ in $L^2$ on the
displayed domain.
It agrees with the pointwise angular formula by the almost sure
interior convergence of $h$.
Projection and finite-energy approximation prove the asserted compression.

Write $\sigma=V\rho_qV^{-1}=\sum_j\lambda_j|F_j\rangle\langle F_j|$.
Lemma~\ref{hfill:lem:collar-moments} implies
\[
 \sum_j\lambda_j\|(E+1)^{3/2}F_j\|^2<\infty.
\]
Equation~\eqref{hfill:eq:multiplication-bound} and Cauchy--Schwarz give
\[
 \int|\Phi_r|D_\sigma\,d\mu
 \leq C_r\left(\operatorname{Tr}(\sigma(E+1)^3)\right)^{1/2}.
\]
Hence summation and integration are justified, and
$\int\Phi_r\,d\nu_\sigma=\operatorname{Tr}(\rho_qW_r)$.

On the energy-$e$ space only $i+j\leq e$ contributes to $W_r$.
It converges there in operator norm to $W$ as $r\uparrow1$.
In the normalized occupation basis, $W$ has nonnegative matrix entries,
and those of $W_r$ lie between zero and the corresponding entries of $W$.
For the vector $|v|$ of absolute values of coefficients, this gives
\[
 |\langle v,W_rv\rangle|
 \leq\langle |v|,W|v|\rangle
 \leq\frac{e(e-1)}2\|v\|^2.
\]
Both operators are symmetric, so the same bound holds for their norms.

The diagonal energy weights of $\rho_q$ have finite second moment.
Dominated summation over the energy spaces proves
\eqref{hfill:eq:radial-limit}.
All higher moments are also finite, so the displayed operator products
have their ordinary trace-class meaning.
For finite-energy states the sum of blocks is finite from the start.
On the energy-two space, $W_r=r^4W$ and $O=2W$,
which gives the last formula.
\end{proof}

The theorem asserts convergence of these expectations.
It does not define an unregularized $\Phi_1$ in $L^2(\mu)$.
The fixed-radius multiplication estimate and the uniform energy-block
bound serve different purposes in its proof.

\subsection{The native and gravitational comparisons}

The supported construction of Section~\ref{hnb:section} connects the
relative coefficient comparison with the scalar filling.
Its domains distinguish the native result from a comparison with gravity.

\begin{remark}[The supported native comparison]\label{hfill:qn:native}
At level one, Section~\ref{hnb:section} constructs the bosonic cochain
complex and the fermionic complex of total charge zero independently.
The cores have finitely many angular weights in each input separately.
Theorem~\ref{hnb:thm:core-fields} gives continuous inverse homotopy
equivalences preserving the vacuum and all ordered neutral current
fields on these cores. Their extensions over the rooted collision complexes
preserve the contact attachments, normal derivatives and cut
restrictions (Proposition~\ref{hnb:prop:supported-comparison}).
On independent-normal charts, the resulting state projection gives
Proposition~\ref{hfill:prop:relative}.

For separated preparations and real collar parameter $0<q\leq1$,
Proposition~\ref{hnb:prop:positive-supported} gives a reflected form
that kills boundaries and is positive definite on its null quotient.
Contacts remain in the supported complex and are null in this form.
Section~\ref{hnb:subsec:collar} constructs the corrected collars
and their higher compatibility homotopies.
When the family includes the normalized annular smearing of
Proposition~\ref{hnb:prop:positive-supported},
Theorem~\ref{hnb:thm:fillings} identifies its physical Hilbert
completion and sewing with the harmonic Gaussian theory.

Translation in Section~\ref{hnb:subsec:translation} is infinitesimal or
coefficientwise formal. A finite translation away from the origin need
not preserve these angular cores. Extension to arbitrary smooth angular
completions or coincident distributional products requires the
corresponding continuous field and multiplication maps.
\end{remark}

\begin{question}\label{hfill:qn:Einstein}
Specify an Einstein-filling state space with boundary conditions,
gauge quotient, action, real form, normalization, insertions and gluing.
Does it admit a structured comparison with the boundary states that
retains the mixed observable of \eqref{hfill:eq:expectations}?
\end{question}

The harmonic scalar field in \eqref{hfill:eq:harmonic} does not supply
an Einstein metric. Nor do these maps identify a half-BPS droplet sector
or an oper state space. These remain distinct geometric
recipients in Question~\ref{hst:qn:gravity}.

